% Introduction — motivate (a) why now, (b) what the catalog is, (c) what
% the search adds, and (d) why the non-overlap-property observation is
% worth a section.

The complexity of multiplying two \(n \times n\) matrices has been
studied since Strassen's 1969 algorithm~\cite{strassen1969} broke the cubic barrier, and
remains open: the exact value of the matrix multiplication exponent
\(\omega\) is unknown, though sub-cubic upper bounds have been driven
down to \(\omega < 2.371\ldots\) via the laser method and its
successors. The asymptotic regime has been productive for theoretical
upper bounds but has produced no concrete algorithms one could implement
for finite \(n\); the constants are astronomical.

Meanwhile, the \emph{small-format} regime -- exact rank bounds for
fixed \(\nmpshape{n}{m}{p}\) tensors with \(n, m, p \le 32\) or so --
has lived a quieter life dominated by hand-constructed schemes:
Strassen (1969)~\cite{strassen1969}, Hopcroft--Kerr
(1971)~\cite{hopcroftkerr1971}, Laderman (1976)~\cite{laderman1976},
Pan (1978, 1980)~\cite{pan1978,pan1980}, and Smirnov (1986,
2013)~\cite{smirnov1986,smirnov2013}. These results are concrete
algorithms; they multiply small matrices using a count of bilinear
products that beats the naïve \(nmp\). When applied recursively, each
small-format rank improvement compounds into an \(\omega\)
improvement. The strongest \emph{practical} small-format exponent comes
from Smirnov's \(\nmpshape{3}{3}{6}\) scheme of rank
40~\cite{smirnov2013}: at \(\omega = 3\log 40 / \log 54 \approx
2.7743\) it sits below the exponent of every explicit cubic scheme in
the catalog, making it the best base case a recursive implementation
can realistically use.

The 2022--2026 period has shaken this landscape, with a fresh record
almost every year. In 2022, AlphaTensor \cite{alphatensor2022} found 47
bilinear products for \(\nmpfield{\Ftwo}{4}{4}{4}\). In 2024,
Kaporin~\cite{kaporin2024brent} gave an explicit complex 48 for
\(\nmpfield{\Complex}{4}{4}{4}\) via a semi-analytical solution of the
Brent equations; in 2025, AlphaEvolve \cite{alphaevolve2025}
independently found 48 for the same shape, which
Dumas--Pernet--Sedoglavic \cite{dps2025} rationalised the same year ---
48 with coefficients in \(\{\pm1,\pm\tfrac12,\pm\tfrac14,\pm\tfrac18\}\),
valid over \(\Rationals\) and hence \(\Reals\). Also in 2025,
Schwartz--Zwecher \cite{schwartzzwecher2025} produced a systematic
disjoint-sum construction giving record ranks at several large shapes. The FMM
catalog at Lille \cite{fmmlille} and Perminov's
\code{FastMatrixMultiplication} repository \cite{perminov} both
aggregate schemes from across the literature \emph{and} contribute
their own results, but via different approaches: FMM layers recursive
derivations (Kronecker with serendipity, concatenation, projection) on top of
known bases, while Perminov runs flip-graph and meta-flip-graph search
to discover fresh low-rank schemes (including ternary-integer ones).

\paragraph{The problem.} These results live in incompatible formats
(Maple scripts, NumPy archives, ad-hoc JSON), use inconsistent shape
conventions (sorted vs unsorted \nmpshape{n}{m}{p}, with or without
explicit field tags), and carry inconsistent attribution. A
\(\nmpshape{3}{3}{3}\) entry sourced from AlphaTensor is sometimes
cited as ``the AlphaTensor scheme'' even though rank 23 was established
by Laderman~\cite{laderman1976} decades earlier -- only some of
AlphaTensor's per-format ranks were genuine improvements over prior
SOTA (its \(\nmpfield{\Ftwo}{4}{4}{4}=47\) among them). The same trap
exists in every bulk-import flow we have audited.

We argue that a catalog that is machine-checkable, fully
attributed, and uniformly serialised is now both possible and
necessary. This paper describes the construction of such a catalog,
together with an exhaustive derivation search that automatically discovers
recombinations of catalog entries.

\paragraph{Field discipline as load-bearing.} A bare \nmpshape{n}{m}{p}
is ambiguous: different fields admit different ranks, as the
\(\nmpshape{4}{4}{4}\) landscape illustrates --- 47 over \Ftwo, 48
over \(\Rationals/\Reals/\Complex\), 49 over \Integers. We keep fields
separate and require every rank claim in the catalog (and in
this paper) to carry a field tag; Section \ref{sec:notation} fixes
the conventions. Commutative algorithms (Waksman
1970~\cite{waksman1970}, Makarov 1986~\cite{makarov1986}, Rosowski
2019~\cite{rosowski2019}) form a separate axis: valid for scalar
matrix multiplication but not liftable to recursive multiplication
over non-commutative rings.

\paragraph{What the catalog catalogs.} Each entry encodes either
\begin{enumerate}[label=(\alph*)]
  \item explicit factor matrices \((U, V, W)\) with field tag,
    addition count, and lineage record; or
  \item a \emph{derived bound} -- a rank given by a closed-form
    construction (Waksman, Rosowski, Pan trilinear aggregation) that we
    reconstruct from its lineage record on demand; or
  \item a \emph{cited bound} -- our last resort, for the edge cases
    where we trust a published rank and attribute it to its source but
    do not (yet) hold its explicit factor matrices.
\end{enumerate}
This three-way split lets us include results we cannot yet
materialise as schemes without misrepresenting their status, and lets
us regenerate the derived layer as our formulas improve.

\paragraph{Derivation over the catalog.} Recursive multiplication
over the same catalog is performed by a search that recombines
existing entries via well-defined operators (Section
\ref{sec:strategies}): axis-flip, Kronecker product (including its
serendipitous ``bud'' variant), axis concatenation, recombination with
allocation (with optional output peeling and pair fusion), and
downward projection -- none of them a disjoint-sum / Sch\"onhage
$\tau$-theorem construction (Section \ref{ssec:disjoint-sum}). The
search is iterated until no further rank improvements appear at any
catalog shape -- a \emph{derivation closure} -- and the resulting wins
are materialised lazily after the search has converged.

\paragraph{The non-overlap property.} An observation that we believe is
underdiscussed: our composition operators never \emph{discover} sharing
between bilinear products while materialising a scheme. A materialised
scheme has exactly the rank the search predicted for it --- the sum
\(\sum_k r_{\text{sub-}k}\) of its sub-schemes' ranks for pure
block-substitution (Kronecker, concatenation, recombination), or a
smaller value where a \emph{known} identity (pair fusion, serendipity,
projection) reduces it, counted before materialisation and never found
during it. What never happens is a product turning out, once expanded,
to do double duty across outputs and merge with another for a free rank
saving. That kind of non-obvious sharing --- as in Strassen's
seven-product \(\nmpfield{\Reals}{2}{2}{2}\) scheme, where a product
serves several output cells at once --- is exactly what hand-crafted
schemes are built around (Laderman, Smirnov, AlphaTensor likewise). We
argue that this makes ``composed scheme matches hand-crafted scheme's
rank'' a meaningful co-existence claim, not a re-discovery claim. The same property
predicts where we \emph{lose}: on the handful of shapes an external
catalog still beats us (six against FMM-Lille), every deficit traces to
one cross-slot sharing device -- a \emph{fusion kernel} -- that our
operators forgo by construction, so those losses confirm the boundary
rather than contradict it. Section \ref{sec:nonoverlap} develops both
directions.

\paragraph{Outline.} Section \ref{sec:notation} fixes notation, the
field-discipline conventions, and the paper's glossary. Section
\ref{sec:architecture} describes the catalog's serialisation, lineage
records, and the explicit/cited/derived split. Section
\ref{sec:strategies} catalogs the derivation operators; Section
\ref{sec:nonoverlap} develops the non-overlap property and its
measured boundary (fusion kernels). Sections \ref{sec:multisets}
and~\ref{sec:hk71} present the two theory contributions ---
recombination multisets, and the constructive realisation of the
Hopcroft--Kerr bound. Section \ref{sec:search} describes the
exhaustive derivation search, and Section \ref{sec:balance} the empirical
case for unbalanced splits. Section \ref{sec:tables} gives the
date-stamped head-to-head against FMM-Lille and Perminov, with an
anatomy of the wins and losses. Section \ref{sec:openquestions}
closes with open questions, and Section \ref{sec:aiagent} records the
human/AI division of labour.
